\begin{figure}[!htbp]
\centering
\resizebox{\linewidth}{!}{%
\begin{tikzpicture}[
    font=\sffamily\footnotesize, >=Latex,
    src/.style={draw=ngBlue, top color=ngBlue!8, bottom color=ngBlue!28,
                rounded corners=3pt, align=center, minimum width=28mm,
                minimum height=11mm},
    proc/.style={draw=ngPurple, top color=ngPurple!8, bottom color=ngPurple!28,
                rounded corners=3pt, align=center, minimum width=28mm,
                minimum height=11mm},
    result/.style={draw=ngGreen, top color=ngGreen!8, bottom color=ngGreen!30,
                rounded corners=3pt, align=center, minimum width=28mm,
                minimum height=11mm},
    impl/.style={draw=ngRed, top color=ngRed!10, bottom color=ngRed!30,
                rounded corners=3pt, align=center, minimum width=36mm,
                minimum height=22mm, font=\footnotesize},
    arr/.style={-{Latex[length=2.5mm]}, line width=0.7pt, draw=ngSlate}
]
% Left column: three raw data sources
\node[src] (s1) at (0,4) {Survey\\(\textit{N}=87, 33 cols)};
\node[src] (s2) at (0,2) {Task ratings\\(\textit{N}=87, 615 cols)};
\node[src] (s3) at (0,0) {App aggregates\\(\textit{N}=310)};

% Second column: validation and linkage processes
\node[proc] (p1) at (3.8,3) {Validation\\artefact removal};
\node[proc] (p2) at (3.8,1) {Linkage by email\\\& anonymisation};

% Third column: derived-variable hub
\node[proc] (p3) at (7.6,2) {Derived variables\\(personal, necessity,\\mention)};

% Fourth column: three RQ results
\node[result] (r1) at (11.4,3.5) {RQ1\,/\,RQ2\\mental models};
\node[result] (r2) at (11.4,2.0) {RQ3\,/\,RQ4\\task necessity};
\node[result] (r3) at (11.4,0.5) {RQ5\\survey$\leftrightarrow$task link};

% Fifth column: implications, placed to the RIGHT of the result column
% so each result connects with a single horizontal-ish arrow and the
% arrows do not cross one another.
\node[impl] (i1) at (15.6,2.0)
    {\textbf{Necessity Gap}\\\textbf{design implications}\\(WWWCF)};

% Source -> processing arrows.
\draw[arr] (s1) -- (p1);
\draw[arr] (s2) -- (p1);
\draw[arr] (s2) -- (p2);
\draw[arr] (s3) -- (p2);

% Processing -> derived-variable hub.
\draw[arr] (p1) -- (p3);
\draw[arr] (p2) -- (p3);

% Hub -> three RQ results.
\draw[arr] (p3) -- (r1);
\draw[arr] (p3) -- (r2);
\draw[arr] (p3) -- (r3);

% Results -> implications.  Three direct arrows from each result box's
% east edge to the implications box's west edge.  Because i1 sits to the
% right of the result column at the same vertical centre as r2, the
% three arrows fan in cleanly without crossing.
\draw[arr] (r1.east) -- (i1.north west);
\draw[arr] (r2.east) -- (i1.west);
\draw[arr] (r3.east) -- (i1.south west);

\end{tikzpicture}%
}
\caption{Analysis architecture. Raw survey and task files (left) feed
the validation and linkage stages, which produce derived variables
that are evaluated against the five research questions. The three RQ
findings converge into the Necessity-Gap design implications (right),
via three direct arrows that fan in from the result column without
crossing any other box.}
\label{fig:analysis-architecture}
\end{figure}
